Large language model (LLM) inference is dominated by repeated dense linear projections whose computational and memory demands increase sharply with model scale. This work studies whether Bit-Sliced Indexing (BSI) can serve as a practical compressed representation for LLM inference on GPUs while retaining a native bitwise execution path rather than falling back to dense low-precision matrix multiplication. The implementation targets PyTorch-based end-to-end inference for OPT-family models and uses a C++/CUDA extension that stores linear weights in BSI form, builds bit-sliced query activations on the fly, and computes matrix products using fused CUDA kernels based on binary tensor-core operations and weighted popcount accumulation. The central question is not whether lower precision is generally faster, but whether a BSI-native execution pipeline can be made substantially more competitive through representation and kernel co-design.

The work grows out of an earlier CPU-side BSI and bitmap foundation in which Bit-Sliced Vectors were built on top of hybrid bitmaps and used to implement dot product, summation, multiplication, and top-k style vector algebra. Those early experiments, including tests on randomly generated vectors and extracted BERT weights, showed that BSI could reduce memory usage and expose useful algebraic structure, but did not yet provide competitive execution time on CPUs. That result motivated the transition to CUDA, where the same BSI semantics were retained while the execution path was redesigned around batched query construction, prebuilt BSI weights, and fused same-bit tensor-core kernels.

The implementation now provides a complete experimental stack: CPU and CUDA BSI builders, a PyTorch replacement for \texttt{torch.nn.Linear}, kernel-only microbenchmarks, end-to-end next-token evaluation on LAMBADA, and comparisons against both dense FP16 baselines and SmoothQuant benchmark scripts. The stable CUDA baseline centers on an SM90/H100 path using \texttt{TM=32}, \texttt{R\_sweep=4}, chunk scaling, and a fixed \((S_a,S_b)=(7,6)\) query/key slice configuration. Under the primary evaluation setting (\texttt{decimal\_places=2}, \texttt{compress\_threshold=0.5}, \texttt{num\_samples=1000}, \texttt{max\_seq\_len=512}, \texttt{fp16} base dtype), the \texttt{facebook/opt-6.7b} model has a dense static size of \texttt{12700.031 MB} and a BSI static size of \texttt{5022.281 MB} under \texttt{scope=all}. In the same setting, BSI achieves \texttt{top1=0.741}, \texttt{top5=0.896}, and \texttt{66.452 ms} dot latency per sample, compared with the dense baseline at \texttt{top1=0.798}, \texttt{top5=0.929}, and \texttt{15.493 ms}.

The main conclusion is that BSI representation alone is insufficient for competitive GPU inference. A generic bit-sliced storage layout still incurs substantial orchestration cost, and the performance gap to dense Torch remains dominated by the interaction between representation, staging, and kernel execution. The contribution of the work is therefore not a claim that BSI already outperforms dense inference, but a complete BSI-native inference stack and a defensible systems result: BSI becomes materially more competitive only when the data layout, chunk handling, and dot kernel are designed together for the target GPU architecture.
